\section{The Model and Its Judgment (RQ3)}
\label{sec:judgment}

The model set the ceiling. Moving from GPT-6.1 Sol to GPT-6.2 Sol changed which tasks passed more
than any harness lever did, and no text we wrote moved the model's choice at a decision point; only
facts put in front of it did.

\para{Most failures are confident deliveries} On the readable tasks, 213 of 227 base failures were
deliveries the model believed done. Each rested on a reading of the request that the model settled
in its first minutes and then verified against itself, so its own checks passed. Side by side with
Codex's public trajectories on the same tasks, Codex settled a stronger design after one long think
and then hardened it, while our runs stopped once their own checks passed. The whole readable gap
to Codex sat in eight tasks that Codex passed reliably (6.75 against 2.1 expected passes), with no
common mechanism.

\subsection{What moved the choice}

We captured the request at a decision point of a real trial and replayed it many times with one
text changed (key-node replay). The replay matched the original prompt exactly, was 98\% cached and
took about 13 seconds a sample. Three kinds of input moved the choice:
\begin{itemize}[leftmargin=1.4em]
  \item \emph{A fact or an order in a user turn at the moment it applies.} A line naming a defect as
    settled, added to the task, made the next move write the unbiased estimator in 23 of 24
    replays. The GO answer of \cref{sec:structure} is of this kind.
  \item \emph{Withholding a misleading cue.} In a planted-choice probe the model kept the option the
    code documents in 16 of 16 replays when told which option that was, and picked the right one in
    16 of 16 when not told.
  \item \emph{Measurements of the data put in front of it.} On one task this took both GPT models
    from 0 of 8 to 8 of 8.
\end{itemize}

\subsection{What did not}

\begin{itemize}[leftmargin=1.4em]
  \item \emph{Advice in the instructions or the profile}, in five wordings. ``A comment is not
    evidence'' moved 0 of 24 decisions, and a profile section written for GPT-6.2 changed no
    measured behaviour.
  \item \emph{Pre-steps by the same model} that list and decide the open points before the work.
    They reached one task only.
  \item \emph{A different model beside the lead as a general reader or chooser.} Claude Opus as a
    reader went 0 of 8 and did harm; as a chooser between two builds its accuracy was 0.55. Another
    model helped only where it knew something GPT did not: in V1 an Opus settler fixed one
    statistics task 4 of 4 times against 0 of 4.
  \item \emph{A note delivered in real time.} An adapter hook placed one general note right after
    the result of the run's first file write, at the same place in every later request, so the
    prompt cache was kept. As a user turn, the note made the next move test the decisive point in
    8 of 8 and 5 to 6 of 8 replays at two captured decisions; as a suffix of the tool result it did
    so in 5 to 6 of 8 and 0 of 8, and without the note in 0 of 4. In whole runs, however, the
    readable 33 scored 21.0 against 21.0 at the same cost: $+1.0$ on the seven open-reading tasks
    it aimed at and $-1.0$ on the other 26. On one task the model ran the requested check, printed
    the score of the form the graders expected (3.239 against its own 3.240), and kept its own
    form.
\end{itemize}

Replays at a captured decision overstate whole runs: a probe-first text took one decision from 0
to 24 of 24 but moved whole runs by $-0.04$ ($\pm 0.85$). A note at the right moment changes the
process, not the judgment, which follows the model's prior. We closed advice as a lever for GPT
judgment on these tasks.

\subsection{Model version and effort}

\begin{table}[t]
\centering
\small
\caption{GPT-6.1 Sol against GPT-6.2 Sol in team 7.0 at max effort on Terminal-Bench~4.0. $S$ is
the sum of per-task mean rewards (\cref{eq:score}); ``best of $N$'' assumes a perfect chooser over
$N$ draws.}
\label{tab:models}
\begin{tabular}{@{}lll@{}}
\toprule
Measure & GPT-6.1 Sol & GPT-6.2 Sol \\
\midrule
Readable 33, $S$ & about 18.8 & about 21 \\
Unread 30, same batch, $S$ & 20.5 & 18.0 (difference $-2.5$, $p = 0.15$) \\
Unread 30, mean of one draw & 19.25 (8 draws) & 17.75 (4 draws) \\
Unread 30, best of 2 / best of 4 & 21.04 / 21.77 & 19.33 / 20.0 \\
Unread tasks that split across draws & 9 of 30 over 8 draws & 6 of 30 over 4 draws \\
Model calls per trial & 46.9 & 31.4 \\
Agent minutes per trial, unread & 33.7 & 28.8 \\
List-price cost per trial, unread & \$1.16 & \$1.15 \\
\bottomrule
\end{tabular}
\end{table}

\para{GPT-6.1 Sol to GPT-6.2 Sol} \Cref{tab:models} compares the two models in the same
configuration. GPT-6.2 made a third fewer model calls and took 15\% fewer minutes at the same
list-price cost per trial. It passes open-reading tasks that GPT-6.1 rarely did (one task 5 of 5
against 0 of 11) and never passes some that GPT-6.1 sometimes did (0 of 8 against 17 of 50 and 16
of 50). It is also more deterministic: its walls are single method choices made the same way in
every draw, which is why sampling more draws helps it less. We kept GPT-6.2 for the final
configuration, since the unread difference was within noise and the agents we compare against run
GPT-6.2.

\para{Effort} In V1, medium effort matched high on ten bug tasks (28/30 against 29/30) at 38\%
less cost, but lost on the multi-bug tasks (0/6 against 3/6, $p \approx 0.06$). On a screen of
readable Terminal-Bench tasks, the earlier lead passed 19 of 38 trials at max, 6 of 21 at medium
and 3 of 15 at high, so high is not a cheap middle. With the lone lead, running at max for the
first 12 calls and at medium afterwards passed 5 of 15 against 11 of 15 at max throughout, at the
same cost. Max effort became the shipped depth for
coding, with the configuration comparison of \cref{sec:structure} as its held-out check.

\para{Answer to RQ3} The harness can change what the model does next (run a check, keep going) but
not what it concludes at a choice; only facts and measurements placed in front of it changed the
choice. The remaining failures are confident single readings, so the levers that can move them must
supply evidence, not advice.
